% Lean module: MonadFormal.Information.Testing.Hellinger.Regularity
\paragraph{Local square-root regularity.}

Consider a real-indexed family of square-root densities $r_\theta$ over
the same $\mu$. Square-root regularity at $\theta$ means differentiability
as an $L^2$-valued function: for some $v\in L^2(\mu;\mathbb R)$,
\[
 r_{\theta+h}=r_\theta+h v+o_{L^2}(|h|).
\]
The norm topology is essential; pointwise differentiation alone does
not control the integrals below. This is the square-root formulation
underlying quadratic-mean regularity in local statistical theory
\cite{duchi2018localTesting}.
\leandef{MonadFormal.Information.HasRootDerivAt}

\begin{theorem}[Local Hellinger and total-variation coefficients]
\label{thm:local-density-distances}
For such a family, write $P_\theta=r_\theta^2\mu$. Then
\[
 \langle r_\theta,v\rangle=0,
 \qquad
 \lim_{h\to0,\,h\ne0}\frac{H^2(P_{\theta+h},P_\theta)}{h^2}
 =\frac12\|v\|_2^2.
\]
Its densities are differentiable in $L^1$, with derivative
$2r_\theta v$, and
\[
 \lim_{h\to0,\,h\ne0}
 \frac{d_{\mathrm{TV}}(P_{\theta+h},P_\theta)}{|h|}
 =\|r_\theta v\|_1.
\]
The coefficients may vanish; nondegeneracy is not automatic.
\end{theorem}
\begin{proof}
Differentiate $\|r_\theta\|_2^2=1$ to obtain
$2\langle r_\theta,v\rangle=0$.
The difference quotient $(r_{\theta+h}-r_\theta)/h$ converges to $v$
in $L^2$. Taking squared norms and dividing by two gives the Hellinger
limit.

Multiplication is a bounded bilinear map
$L^2\times L^2\to L^1$, since
$\|fg\|_1\le\|f\|_2\|g\|_2$ by Cauchy--Schwarz.
Its product rule therefore gives
\[
 \frac{r_{\theta+h}^2-r_\theta^2}{h}
 \longrightarrow 2r_\theta v\quad\text{in }L^1.
\]
Taking norms is continuous, and \Cref{thm:hellinger-tv} identifies
total variation with half the $L^1$ norm of the density difference.
Thus its quotient by $|h|$ tends to
$\tfrac12\|2r_\theta v\|_1=\|r_\theta v\|_1$.
\leanproof{MonadFormal.Information.HasRootDerivAt.orthogonal,MonadFormal.Information.HasRootDerivAt.hellinger_limit,MonadFormal.Information.l2Product,MonadFormal.Information.l2Product_coe,MonadFormal.Information.l2Product_norm_le,MonadFormal.Information.l2Product_comm,MonadFormal.Information.l2ProductCLM,MonadFormal.Information.RootDensity.density,MonadFormal.Information.RootDensity.density_coe,MonadFormal.Information.RootDensity.tv_eq_density_norm,MonadFormal.Information.HasRootDerivAt.density_deriv,MonadFormal.Information.HasRootDerivAt.tv_limit}
\end{proof}

% monad-open: information-profile
These are single-observation, scalar-parameter expansions. Passing
to independent product laws or a limiting optimal test is a further
step. A KL expansion is different again: the argument above never
differentiates a log likelihood and supplies no bound on its tails
or on likelihood ratios near zero.
